% !TeX root = ../教材题库_全部条目_转写工作稿.tex
% Source: XJ-B1-C06-HW-P214-Y1026
% 第6章 / 6.1 获取数据的途径与统计的一些基本概念 / 教材 p.214 / PDF 第221页
% OCR未校对：公式与文字均可能有误；未完成数学TeX转写；禁止直接发给学生
% 原题及相邻原解截图（核对用）：../题目原图/XJ-B1-C06-HW-P214-Y1026-01.webp
% 本文件只在转写完成后作为题干使用；OCR 条目尚不是可靠题干。
\begin{bankitem}{习题6.1 第5题}{第6章 / 6.1 获取数据的途径与统计的一些基本概念 / 教材 p.214 / PDF 第221页}{XJ-B1-C06-HW-P214-Y1026}{OCR未校对：公式与文字均可能有误；未完成数学TeX转写；禁止直接发给学生}
\begin{ocrdraft}
[教材p.214 / PDF第221页]
\par
5．数据的真实性．
\par
当下，许多媒体充斥着各种各样的统计数字和图表．“让数据说话”成为许多
\par
广告的常用手法．例如，某减肥药厂商做广告时声称，其生产的减肥药有效率达到
\par
80％．见到这样的广告数据你会怎么想？试与同学从各种媒体中收集一些广告，并
\par
用统计知识分析一下他们所提供的数据和结论的真实性．（提示：数据是谁收集的？
\par
怎样收集的？样本容量是多少？样本是如何选取的？这些数据是什么时候收集的？）
\end{ocrdraft}
\end{bankitem}
